\section{Related Work}\label{sec:related}

\textit{Tracking warnings across versions.} Spacco, Hovemeyer and Pugh treated warning identity as a problem of its own. They added two tracking techniques to FindBugs, so that reviewers could remember warnings already judged safe and see which ones were new, and applied them to Sun's Java runtime library~\cite{spacco2006}. Later work used tracking as a means to another end: Avgustinov et al.\ followed violations through revision histories at Semmle to credit each developer with the violations they introduced or removed~\cite{avgustinov2015}, and Kim and Ernst mined warning histories to rank warning categories by how often fixes removed them~\cite{kim2007}. In both, the tracker is part of the method, and its errors reach the results unmeasured.

The closest work is by Li and Yang~\cite{li2022statictracker}. They labelled the evolution of 3,451 SpotBugs and PMD warnings in four Java projects by hand, showed that an existing tracker failed on refactorings, code shifts and volatile metadata, and proposed StaticTracker, which raised tracking precision from 64.4\% to 90.3\%. Our question is different. We do not build a better tracker; we compare the identity methods that shipped tools use (SonarQube, Bandit, Semgrep and GitHub code scanning), each ported from its source and checked against the tool, on the same findings from 35 Python and JavaScript applications. We also count the two kinds of error separately. A baseline that reports old findings as new is a nuisance; one that hides new findings is a security problem.

\textit{Why identity matters in practice.} Developers name false positives and the way warnings are presented as barriers to using static analysis~\cite{johnson2013}, and at Google, the analyses developers act on show their findings in code review, before the code is committed~\cite{sadowski2015tricorder,sadowski2018lessons}. Showing a developer only the findings their change introduced needs exactly the identity we measure. Suppressions depend on it too: Hu et al.\ found that half of the suppression comments in their study suppress no warning at all~\cite{hu2025suppressed}. SARIF~2.1.0 gives tools fields for fingerprints and for a result's baseline state~\cite{sarif2020}, yet the tools we compare match results in very different ways (Section~\ref{sec:schemes}).
